\section{AHSME 1996}
        \begin{enumerate}
        	\item (\textit{AHSME 1996}) The addition below is incorrect. What is the largest digit that can be changed to make the addition correct?


 $\begin{tabular}{rr}&\ \texttt{6 4 1}\\ &\texttt{8 5 2}\\ +&\texttt{9 7 3}\\ \hline  &\texttt{2 4 5 6}\end{tabular}$
$\text{(A)}\ 4\qquad\text{(B)}\ 5\qquad\text{(C)}\ 6\qquad\text{(D)}\ 7\qquad\text{(E)}\ 8$

 \textbf{Answer: }D



\textbf{Solution}

Doing the addition as is, we get $641 + 852 + 973 = 2466$.  This number is $10$ larger than the desired sum of $2456$.  Therefore, we must make one of the three numbers $10$ smaller.

 We may either change $641 \rightarrow 631$, $852 \rightarrow 842$, or $973 \rightarrow 963$.  Either change results in a valid sum.  The largest digit that could be changed is thus the $7$ in the number $973$, and the answer is $\boxed{\textbf{(D) }7}$.

 


	\item (\textit{AHSME 1996}) Each day Walter gets $3$ dollars for doing his chores or $5$ dollars for doing them exceptionally well. After $10$ days of doing his chores daily, Walter has received a total of $36$ dollars. On how many days did Walter do them exceptionally well?


 $\text{(A)}\ 3\qquad\text{(B)}\ 4\qquad\text{(C)}\ 5\qquad\text{(D)}\ 6\qquad\text{(E)}\ 7$

 \textbf{Answer: }A



\textbf{Solution 1}

If Walter had done his chores for $10$ days without doing any of them well, he would have earned $3 \cdot 10 = 30$ dollars.  He got $6$ dollars more than this.

 He gets a $5 - 3 = 2$ dollar bonus every day he does his chores well.  Thus, he did his chores exceptionally well $\frac{6}{2} = 3$ days, and the answer is $\boxed{A}$.

 



\textbf{Solution 2}

If Walter had done his chores for $10$ days exceptionally well, he would have earned $5 \cdot 10 = 50$ dollars.  He got $50 - 36 = 14$ dollars less than this.

 He gets $2$ dollars docked from his pay if he doesn't do his chores well.  Therefore, he didn't do his chores well on $\frac{14}{2} = 7$ days.  The other $10 - 7 = 3$ days, he did them exceptionally well.  Therefore, the answer is $\boxed{A}$.

 



\textbf{Solution 3}

Let $b$ be the number of days Walter does his chores but doesn't do them well, and let $w$ be the number of days he does his chores exceptionally well.

 $b + w = 10$ since there are $10$ days Walter does chores.

 $3b + 5w = 36$ since $3b$ is the amount he earns from doing his chores not well, and $5w$ is the amount he earnes from doing his chores exceptionally well, and those two sum to $36$ dollars.

 Multiply the first equation by $3$ to get:

 $3b + 3w = 30$
$3b + 5w = 36$

 Subtract the first equation from the second equation to get:

 $5w - 3w = 36 - 30$
$2w = 6$
$w = 3$

 Thus, he does his chores well $3$ days, and the answer is $\boxed{A}$.

 


	\item (\textit{AHSME 1996}) $\frac{(3!)!}{3!}=$
$\text{(A)}\ 1\qquad\text{(B)}\ 2\qquad\text{(C)}\ 6\qquad\text{(D)}\ 40\qquad\text{(E)}\ 120$

 \textbf{Answer: }E



\textbf{Solution 1}


 \begin{align*} \frac{6!}{6} &= \frac{\cancel{6} \cdot 5!}{\cancel{6}} \\ &= 5! = \boxed{\textbf{(E) }120} \end{align*}





\textbf{Solution 2}


 \begin{align*} \frac{(3!)!}{3!} &= \frac{6!}{3!} \\ &= \frac{720}{6} \\ &= \boxed{\textbf{(E) }120} \end{align*}




	\item (\textit{AHSME 1996}) Six numbers from a list of nine integers are $7,8,3,5, 9$ and $5$. The largest possible value of the median of all nine numbers in this list is


 $\text{(A)}\ 5\qquad\text{(B)}\ 6\qquad\text{(C)}\ 7\qquad\text{(D)}\ 8\qquad\text{(E)}\ 9$

 \textbf{Answer: }D



\textbf{Solution}

First, put the six numbers we have in order, since we are concerned with the median:  $3, 5, 5, 7, 8, 9$.

 We have three more numbers to insert into the list, and the median will be the $5^{th}$ highest (and $5^{th}$ lowest) number on the list.  If we top-load the list by making all three of the numbers greater than $9$, the median will be the highest it can possibly be.  Thus, the maximum median is the fifth piece of data in the list, which is $8$, giving an answer of $\boxed{D}$.

 (In fact, as long as the three new integers are greater than $8$, the median will be $8$.)

 This illustrates one important fact about medians:  no matter how high the three "outlier" numbers are, the median will never be greater than $8$.  The arithmetic mean is, generally speaking, more sensitive to such outliers, while the median is resistant to a small number of data that are either too high or too low.  

 


	\item (\textit{AHSME 1996}) Given that $0 < a < b < c < d$, which of the following is the largest?


 $\text{(A)}\  \frac{a+b}{c+d} \qquad\text{(B)}\ \frac{a+d}{b+c} \qquad\text{(C)}\  \frac{b+c}{a+d} \qquad\text{(D)}\  \frac{b+d}{a+c} \qquad\text{(E)}\ \frac{c+d}{a+b}$

 \textbf{Answer: }E



\textbf{Solution 1}

Assuming that one of the above fractions is indeed always the largest, try plugging in $a=1, b=2, c=3, d=4$, since those are valid values for the variables given the constraints of the problem.  The options become:

 $\text{(A)}\  \frac{1+2}{3+4} \qquad\text{(B)}\ \frac{1+4}{2+3} \qquad\text{(C)}\  \frac{2+3}{1+4} \qquad\text{(D)}\  \frac{2+4}{1+3} \qquad\text{(E)}\ \frac{3+4}{1+2}$

 Simplified, the options are $\frac{3}{7}, 1, 1, \frac{3}{2}, \frac{7}{3}$, respectively.  Since $\frac{7}{3}$ is the only option that is greater than $2$, the answer is $\boxed{E}$.

 



\textbf{Solution 2}

To make a fraction large, you want the largest possible numerator, and the smallest possible denominator.  Since $c$ and $d$ are the two largest numbers, they should go in the numerator as a sum.  Since $a$ and $b$ are the smallest numbers, they should go in the denominator as a sum.  Thus, the answer is $\boxed{E}$.

 You can compare option $E$ with every other fraction:  all numerators are smaller than $E$'s numerator, and all denominators are larger than $E$'s denominator.

 


	\item (\textit{AHSME 1996}) If $f(x) = x^{(x+1)}(x+2)^{(x+3)}$, then $f(0)+f(-1)+f(-2)+f(-3) =$
$\text{(A)}\ -\frac{8}{9}\qquad\text{(B)}\ 0\qquad\text{(C)}\ \frac{8}{9}\qquad\text{(D)}\ 1\qquad\text{(E)}\ \frac{10}{9}$

 \textbf{Answer: }E



\textbf{Solution}

Plugging in $x=0$ into the function will give $0^1\cdot 2^3$.  Since $0^1 = 0$, this gives $0$.

 Plugging in $x=-1$ into the function will give $(-1)^0 \cdot 1^2$.  Since $(-1)^0 = 1$ and $1^2 = 1$, this gives $1$.

 Plugging in $x=-2$ will give a $0^1$ factor as the second term, giving an answer of $0$.

 Plugging in $x=-3$ will give $(-3)^{-2}\cdot (-1)^0$.  The last term is $1$, while the first term is $\frac{1}{(-3)^2} = \frac{1}{9}$

 Adding up all four values, the answer is $1 + \frac{1}{9} = \frac{10}{9}$, and the right answer is $\boxed{E}$.

 


	\item (\textit{AHSME 1996}) A father takes his twins and a younger child out to dinner on the twins' birthday. The restaurant charges $4.95$ for the father and $0.45$ for each year of a child's age, where age is defined as the age at the most recent birthday. If the bill is $9.45$, which of the following could be the age of the youngest child? 


 $\text{(A)}\ 1\qquad\text{(B)}\ 2\qquad\text{(C)}\ 3\qquad\text{(D)}\ 4\qquad\text{(E)}\ 5$

 \textbf{Answer: }B



\textbf{Solution}

The bill for the three children is $9.45 - 4.95 = 4.50$.  Since the charge is $0.45$ per year for the children, they must have $\frac{4.50}{0.45} = 10$ years among the three of them.

 The twins must have an even number of years in total (presuming that they did not dine in the 17 minutes between the time when the first twin was born and the second twin was born).  If we let the twins be $5$ years old, that leaves $10 - 2\cdot 5 = 0$ years leftover for the youngest child.  But if the twins are each $4$ years old, then the youngest child could be $10 - 2\cdot 4 = 2$ years old, which is choice $\boxed{B}$.

 Of note, this is the only possibility for the ages of the children.  If the twins were $3$ years old, the "younger" child would be $10 - 2\cdot 3 = 4$ years old.

 


	\item (\textit{AHSME 1996}) If $3 = k\cdot 2^r$ and $15 = k\cdot 4^r$, then $r =$
$\text{(A)}\ -\log_{2}5\qquad\text{(B)}\ \log_{5}2\qquad\text{(C)}\ \log_{10}5\qquad\text{(D)}\ \log_{2}5\qquad\text{(E)}\ \frac{5}{2}$

 \textbf{Answer: }D



\textbf{Solution}

We want to find $r$, so our strategy is to eliminate $k$.

 The first equation gives $k = \frac{3}{2^r}$.

 The second equation gives $k = \frac{15}{4^r}$

 Setting those two equal gives $\frac{3}{2^r} = \frac{15}{4^r}$

 Cross-multiplying and dividing by $3$ gives $5\cdot 2^r = 4^r$.

 We know that $4^r = 2^r \cdot 2^r$, so we can divide out $2^r$ from both sides (which is legal since $2^r \neq 0$), and we get:

 $5 = 2^r$
$r = \log_2 5$, which is option $\boxed{D}$.

 



\textbf{Solution 2}

Take corresponding logs and split up each equation to obtain:

 $\log_2 3 = (\log_2 k) + (r)$
$\log_4 15 = (log_4 k) + (r)$

 Then subtract the log from each side to isolate r:

 $\log_2 (\frac{3}{k}) = r$
$\log_4 (\frac{15}{k}) = r$

 Then set equalities and solve for k:

 $\log_2 (\frac{3}{k}) = \log_4 (\frac{15}{k})$
$\frac{3}{k} = \sqrt{\frac{15}{k}}$

 After solving we find that $k = \frac{3}{5}$. Plugging into either of the equations and solving (easiest with equation 1) we find that $r = \log_2 5 \implies \boxed{D}$




\textbf{Solution 3}

Since we have a system we can divide both equations: $\frac{15=k \cdot 4^r}{3=k \cdot 2^r}$. It is also important to know that $4^r$ is the same as $2^{2r}$. Simplifying the division of both equations you get $5=2^r$. Using the basic log formula $\log_m n = x \implies m^x=n$. Since $m^x=n$ plugging the numbers in makes $m=2$, $x=r$, $n=5$. This means $r$ is $\boxed{\textbf{(D)}\log_2 5}$



	\item (\textit{AHSME 1996}) Triangle $PAB$ and square $ABCD$ are in perpendicular planes. Given that $PA = 3, PB = 4$ and $AB = 5$, what is $PD$? 

                    \begin{figure}[H]
                        \centering
                        \includegraphics[width=0.5\textwidth]{Problem/AHSME_I_1996_9_Problem_Figure_1.png}
                    \end{figure}
                    
$\text{(A)}\ 5\qquad\text{(B)}\ \sqrt{34} \qquad\text{(C)}\ \sqrt{41}\qquad\text{(D)}\ 2\sqrt{13}\qquad\text{(E)}\ 8$

 \textbf{Answer: }B



\textbf{Solution 1}

Since the two planes are perpendicular, it follows that $\triangle PAD$ is a \href{https://artofproblemsolving.com/wiki/index.php?title=Right\_triangle}{right triangle}. Thus, $PD = \sqrt{PA^2 + AD^2} = \sqrt{PA^2 + AB^2} = \sqrt{34}$, which is option $\boxed{\text{B}}$.

 



\textbf{Solution 2}

Place the points on a coordinate grid, and let the $xy$ plane (where $z=0$) contain triangle $PAB$.  Square $ABCD$ will have sides that are vertical.

 Place point $P$ at $(0,0,0)$, and place $PA$ on the x-axis so that $A(3,0,0)$, and thus $PA = 3$.

 Place $PB$ on the y-axis so that $B(0,4,0)$, and thus $PB = 4$.  This makes $AB = 5$, as it is the hypotenuse of a 3-4-5 right triangle (with the right angle being formed by the x and y axes).  This is a clean use of the fact that $\triangle PAB$ is a right triangle.

 Since $AB$ is one side of $\square ABCD$ with length $5$, $BC = 5$ as well.  Since $BC \perp AB$, and $BC$ is also perpendicular to the $xy$ plane, $BC$ must run stright up and down.  WLOG pick the up direction, and since $BC = 5$, we travel $5$ units up to $C(0,4,5)$.  Similarly, we travel $5$ units up from $A(3,0,0)$ to reach $D(3,0,5)$.

 We now have coordinates for $P$ and $D$.  The distance is $\sqrt{(5-0)^2 + (0-0)^2 + (3-0)^2} = \sqrt{34}$, which is option $\boxed{\text{B}}$.

 


	\item (\textit{AHSME 1996}) How many line segments have both their endpoints located at the vertices of a given cube?


 $\text{(A)}\ 12\qquad\text{(B)}\ 15\qquad\text{(C)}\ 24\qquad\text{(D)}\ 28\qquad\text{(E)}\ 56$

 \textbf{Answer: }D



\textbf{Solution 1}

There are $8$ choices for the first endpoint of the line segment, and $7$ choices for the second endpoint, giving a total of $8\cdot 7 = 56$ segments.  However, both $\overline{AB}$ and $\overline{BA}$ were counted, while they really are the same line segment.  Every segment got double counted in a similar manner, so there are really $\frac{56}{2} = 28$ line segments, and the answer is $\boxed{D}$.

 In shorthand notation, we're choosing $2$ endpoints from a set of $8$ endpoints, and the answer is $\binom{8}{2} = \frac{8!}{6!2!} = 28$.

 



\textbf{Solution 2}

Each segment is either an edge, a facial diagonal, or a long/main/spacial diagonal.

 A cube has $12$ edges:  Four on the top face, four on the bottom face, and four that connect the top face to the bottom face.

 A cube has $6$ square faces, each of which has $2$ facial diagonals, for a total of $6\cdot 2 = 12$.

 A cube has $4$ spacial diagonals:  each diagonal goes from one of the bottom vertices to the "opposite" top vertex.

 Thus, there are $12 + 12 + 4 = 28$ segments, and the answer is $\boxed{D}$.

 


	\item (\textit{AHSME 1996}) Given a circle of radius $2$, there are many line segments of length $2$ that are tangent to the circle at their midpoints. Find the area of the region consisting of all such line segments.


 $\text{(A)}\ \frac{\pi} 4\qquad\text{(B)}\ 4-\pi\qquad\text{(C)}\ \frac{\pi} 2\qquad\text{(D)}\ \pi\qquad\text{(E)}\ 2\pi$

 \textbf{Answer: }D



\textbf{Solution}

Let line segment $AB = 2$, and let it be tangent to circle $O$ at point $P$, with radius $OP = 2$.  Let $AP = PB = 1$, so that $P$ is the midpoint of $AB$.

 $\triangle OAP$ is a right triangle with right angle at $P$, because $AB$ is tangent to circle $O$ at point $P$, and $OP$ is a radius.

 Since $AP^2 + OP^2 = OA^2$ by the \href{https://artofproblemsolving.com/wiki/index.php?title=Pythagorean\_Theorem}{Pythagorean Theorem}, we can find that $OA = \sqrt{1^2 + 2^2} = \sqrt{5}$. Similarly, $OB = \sqrt{5}$ also.

 Line segment $APB$ can rotate around the circle.  The closest distance of this segment to the center will always be $OP = 2$, and the longest distance of this segment will always be $PA = PB = \sqrt{5}$.  Thus, the region in question is the \href{https://artofproblemsolving.com/wiki/index.php?title=Annulus}{annulus} of a circle with outer radius $\sqrt{5}$ and inner radius $2$.  This area is $\pi \cdot 5 - \pi \cdot 4 = \pi$, and the answer is $\boxed{D}$.

 


	\item (\textit{AHSME 1996}) A function $f$ from the integers to the integers is defined as follows: 


 
 \[f(x) =\begin{cases}n+3 &\text{if n is odd}\\ n/2 &\text{if n is even}\end{cases}\]


 Suppose $k$ is odd and $f(f(f(k))) = 27$.  What is the sum of the digits of $k$?


 $\textbf{(A)}\ 3\qquad\textbf{(B)}\ 6\qquad\textbf{(C)}\ 9\qquad\textbf{(D)}\ 12\qquad\textbf{(E)}\ 15$

 \textbf{Answer: }B



\textbf{Solution}

\textbf{First iteration} To get $f(k) = 27$, you could either have $f(27 - 3)$ and add $3$, or $f(27\cdot 2)$ and divide by $2$.

 If you had the former, you would have $f(24)$, and the function's rule would have you divide.  Thus, $k=54$ is the only number for which $f(k) = 27$.

 \textbf{Second iteration} Going out one step, if you have $f(f(k)) = 27$, you would have to have $f(k) = 54$.   For $f(k) = 54$, you would either have $f(54-3)$ and add $3$, or $f(54\cdot 2)$ and divide by $2$.

 Both are possible:  $f(51)$ and $f(108)$ return values of $54$.  Thus, $f(f(51)) = f(54) = 27$, and $f(f(108)) = f(54) = 27$.

 


 \textbf{Third iteration} Going out the final step, if you have $f(f(f(k))) = 27$, you would have to have $f(f(k))) = 51$ or $f(f(k)) = 108$.

 If you doubled either of these, $k$ would not be odd.  So you must subtract $3$.

 If you subtract $3$ from $51$, you would compute $f(48)$, which would halve it, and not add the $3$ back.

 If you subtract $3$ from $108$, you would compute $f(105)$, which would add the $3$ back.

 Thus, $f(f(f(105))) = f(f(108)) = f(54) = 27$, and $105$ is odd.  The desired sum of the digits is $6$, and the answer is $\boxed{B}$.

 



\textbf{Solution 2 (rigorous but easy)}


 \[f(f(f(k))) = 27\]

 We will work from the inside to the outside and alternate \( k \) between even and odd.

 If \( k \) is odd, then \( f(k) \) in terms of \( k \) is \(\frac{k+3}{2}\).

 Next, we have \( f\left(\frac{k+3}{2}\right) \), and \(\frac{k+3}{2}\) is even, so \( f\left(\frac{k+3}{2}\right) = \frac{k+3}{4} \).

 We are given that \( f\left(\frac{k+3}{4}\right) = 27 \).

 
 \[\frac{k+3}{4} = 27 \implies k = 105.\]

 The sum of its digits is \( 1 + 0 + 5 = 6 \).

 


 {gnv12}

 


	\item (\textit{AHSME 1996}) Sunny runs at a steady rate, and Moonbeam runs $m$ times as fast, where $m$ is a number greater than 1.  If Moonbeam gives Sunny a head start of $h$ meters, how many meters must Moonbeam run to overtake Sunny? 


 $\text{(A)}\ hm\qquad\text{(B)}\ \frac{h}{h+m}\qquad\text{(C)}\ \frac{h}{m-1}\qquad\text{(D)}\ \frac{hm}{m-1}\qquad\text{(E)}\ \frac{h+m}{m-1}$

 \textbf{Answer: }D



\textbf{Solution}

If Sunny runs at a rate of $s$ for $h$. Then the distance covered is $sh$. Now we know that Moonbeam runs $m$ times as fast than Sunny, so Moonbeam runs at the rate of $ms$. Now Moonbeam gave Sunny a headstart of $h$ meters, so he will catch  on Sunny at the rate of $s(m-1)$ . At time $\frac{h}{m-1}$  Moon beam will catch on Sunny. Now  we are asked how much in meters he have to run to catch on Sunny. That is $\frac{hm}{m-1}$.

 



\textbf{Solution 2}

Note that $h$ is a length, while $m$ is a dimensionless constant.  Thus, $h$ and $m$ cannot be added, and $B$ and $E$ are not proper answers, since they both contain $h+m$.

 Thus, we only concern ourselves with answers $A, C, D$.

 If $m$ is a very, very large number, then Moonbeam will have to run just over $h$ meters to reach Sunny.  Or, in the language of limits:

 $\lim_{m\rightarrow \infty} d(m) = h$, where $d(m)$ is the distance Moonbeam needs to catch Sunny at the given rate ratio of $m$.

 In option $A$, when $m$ gets large, the distance gets large.  Thus, $A$ is not a valid answer.

 In option $C$, when $m$ gets large, the distance approaches $0$, not $h$ as desired.  This is not a valid answer.  (In fact, this is the distance Sunny runs, which does approach $0$ as Moonbeam gets faster and faster.)

 In option $D$, when $m$ gets large, the ratio $\frac{m}{m-1}$ gets very close to, but remains just a tiny bit over, the number $1$.  Thus, when you multiply it by $h$, the ratio in option $D$ gets very close to, but remains just a tiny bit over, $h$.  Thus, the best option out of all the choices is $\boxed{D}$.

 



\textbf{Solution 3}

Assume that Sunny originally runs at a unit speed, and thus Moonbeam runs at a rate of $m$.

 Choose a new reference frame where Sunny is still, and Moonbeam runs at a rate of $m-1$.  In this new reference frame, the distance to be run is still $h$.

 Moonbeam runs this distance $h$ in a time of $\frac{h}{m-1}$

 Returning to the original reference frame, if Moonbeam runs for $\frac{h}{m-1}$ seconds, Moonbeam will cover a distance of $\frac{hm}{m-1}$, which is option $\boxed{D}$.

 


	\item (\textit{AHSME 1996}) Let $E(n)$ denote the sum of the even digits of $n$. For example, $E(5681) = 6+8 = 14$.  Find $E(1)+E(2)+E(3)+\cdots+E(100)$
$\text{(A)}\ 200\qquad\text{(B)}\ 360\qquad\text{(C)}\ 400\qquad\text{(D)}\ 900\qquad\text{(E)}\ 2250$

 \textbf{Answer: }C



\textbf{Solution}

The problem is asking for the sum of all the even digits in the numbers $1$ to $100$.  We can remove $100$ from the list, add $00$ to the list, and tack on some leading zeros to the single digit numbers without changing the sum of the even digits.  This gives the list:

 $00, 01, 02, 03, ..., 10, 11, ..., 98, 99$.

 There are $2\cdot 100 = 200$ digits on that list, and each digit appears $\frac{200}{10} = 20$ times.

 Thus, each even digit appears $20$ times, and the sum of all the even digits is $0 \cdot 20 + 2\cdot 20 + 4\cdot 20 + 6\cdot 20 + 8\cdot 20 = (0 + 2 + 4 + 6 + 8)\cdot 20 = 400$, and the correct answer is $\boxed{C}$.

 


	\item (\textit{AHSME 1996}) Two opposite sides of a rectangle are each divided into $n$ congruent segments, and the endpoints of one segment are joined to the center to form triangle $A$. The other sides are each divided into $m$ congruent segments, and the endpoints of one of these segments are joined to the center to form triangle $B$. [See figure for $n=5, m=7$.] What is the ratio of the area of triangle $A$ to the area of triangle $B$?


 
                    \begin{figure}[H]
                        \centering
                        \includegraphics[width=0.5\textwidth]{Problem/AHSME_I_1996_15_Problem_Figure_1.png}
                    \end{figure}
                    
$\text{(A)}\ 1\qquad\text{(B)}\ m/n\qquad\text{(C)}\ n/m\qquad\text{(D)}\ 2m/n\qquad\text{(E)}\ 2n/m$

 \textbf{Answer: }B



\textbf{Solution 1}

Place the rectangle on a coordinate grid, with diagonal vertices $(0,0)$ and $(x,y)$.  Each horizontal segment of the rectangle will have length $\frac{x}{m}$, while each vertical segment of the rectangle will have length $\frac{y}{n}$.

 The center of this rectangle will be $(\frac{x}{2}, \frac{y}{2})$.

 Triangle $A$ has a base length of $\frac{y}{n}$, one of the vertical segments.  It has an altitude of $\frac{x}{2}$, which is the perpendicular distance from the center of the square to the left side of the square.  Thus, the area of triangle $A$ is $\frac{1}{2}\cdot\frac{y}{n}\cdot\frac{x}{2} = \frac{xy}{4n}$.

 Triangle $B$ has a base length of $\frac{x}{m}$, one of the horizontal segments.  It has an altitude of $\frac{y}{2}$, which is the perpendicular distance from the center of the square to the bottom side of the square.  Thus, the area of triangle $B$ is $\frac{1}{2}\cdot\frac{x}{m}\cdot\frac{y}{2} = \frac{xy}{4m}$.

 The ratio of areas is  $\frac{\frac{xy}{4n}}{\frac{xy}{4m}} = \frac{m}{n}$, which is answer $\boxed{B}$.

 



\textbf{Solution 2}

Alternately, the algebra can be made simpler if you arbitrarily assume that $(x,y) = (m,n)$, and thus that each side of the rectangle is cut into unit segments.  In that case, the ratio of the areas of the two triangles $A$ and $B$ that have the same base length is just the ratio of the heights.  Triangle $A$ would have height $\frac{m}{2}$, while triangle $B$, while triangle $B$ would have height $\frac{n}{2}$, giving a ratio of $\frac{m}{n}$, which is answer $\boxed{B}$.

 This solution makes the extra assumption that the rectangle has dimensions $m$ by $n$, instead of arbitrary $x$ by $y$ dimensions, and is not a formal "proof"; but since the answer is invariant, extra assumptions will not change the solution.

 


	\item (\textit{AHSME 1996}) A fair standard six-sided dice is tossed three times. Given that the sum of the first two tosses equal the third, what is the probability that at least one "2" is tossed?


 $\text{(A)}\ \frac{1}{6}\qquad\text{(B)}\ \frac{91}{216}\qquad\text{(C)}\ \frac{1}{2}\qquad\text{(D)}\ \frac{8}{15}\qquad\text{(E)}\ \frac{7}{12}$

 \textbf{Answer: }D



\textbf{Solution}

The third toss cannot be $1$, since the minimal sum on the other two tosses is $2$.

 If the third roll is $2$, then the first two rolls must be $(1,1)$.

 If the third roll is $3$, then the first two rolls must be $(1,2)$ or $(2,1)$.

 If the third roll is $4$, then the first two rolls must be $(1,3)$, $(2,2)$, or $(3,1)$.

 If the third roll is $5$, then the first two rolls must be $(1,4)$, $(2,3)$, $(3,2)$, or $(4,1)$.

 If the third roll is $6$, then the first two rolls must be $(1,5)$, $(2,4)$, $(3,3)$, $(4,2)$, or $(5,1)$

 There are $15$ possibilities for the three tosses.  Of those possibilities, $7$ have a $2$ in the first two rolls, and $1$ has a $2$ in the third.  Therefore, the answer is $\boxed{\frac{8}{15}}$.

 


	\item (\textit{AHSME 1996}) In rectangle $ABCD$, angle $C$ is trisected by $\overline{CF}$ and $\overline{CE}$, where $E$ is on $\overline{AB}$, $F$ is on $\overline{AD}$, $BE=6$ and $AF=2$. Which of the following is closest to the area of the rectangle $ABCD$?

                    \begin{figure}[H]
                        \centering
                        \includegraphics[width=0.5\textwidth]{Problem/AHSME_I_1996_17_Problem_Figure_1.png}
                    \end{figure}
                    
$\text{(A)}\ 110\qquad\text{(B)}\ 120\qquad\text{(C)}\ 130\qquad\text{(D)}\ 140\qquad\text{(E)}\ 150$

 \textbf{Answer: }E



\textbf{Solution}

Since $\angle C = 90^\circ$, each of the three smaller angles is $30^\circ$, and $\triangle BEC$ and $\triangle CDF$ are both $30-60-90$ triangles.

 
                    \begin{figure}[H]
                        \centering
                        \includegraphics[width=0.5\textwidth]{Solution/AHSME_I_1996_17_Solution_Figure_1.png}
                    \end{figure}
                    

 Defining the variables as illustrated above, we have $x = 6\sqrt{3}$ from $\triangle BEC$

 Then $x-2 = 6\sqrt{3} - 2$, and $y = \sqrt{3} (6 \sqrt{3} - 2) = 18 - 2\sqrt{3}$.

 The area of the rectangle is thus $xy = 6\sqrt{3}(18 - 2\sqrt{3}) = 108\sqrt{3} - 36$.

 Using the approximation $\sqrt{3} \approx 1.7$, we get an area of just under $147.6$, which is closest to answer $\boxed{\text{E}}$.  (The actual area is actually greater, since $\sqrt{3} > 1.7$).

 



\textbf{Solution 1.1 (Better)}

Use the process above, but use $\sqrt{3} \approx 1.73$. You should get $[ABCD]=150.84$, which then you select $\boxed{\text{E}}$. Notice that the actual area, when plugged into a calculator, yields about $151.0614872$.

 


	\item (\textit{AHSME 1996}) A circle of radius $2$ has center at $(2,0)$. A circle of radius $1$ has center at $(5,0)$. A line is tangent to the two circles at points in the first quadrant. Which of the following is closest to the $y$-intercept of the line?


 $\text{(A)}\ \sqrt{2}/4\qquad\text{(B)}\ 8/3\qquad\text{(C)}\ 1+\sqrt 3\qquad\text{(D)}\ 2\sqrt 2\qquad\text{(E)}\ 3$

 \textbf{Answer: }D



\textbf{Solution}

The two circles are tangent to each other at the point $(4,0)$, since it is both $2$ units from $(2,0)$ and $1$ unit from $(5,0)$.

 Label the x-intercept of the common tangent line $A$, and label the y-intercept of the common tangent $B$.  Triangle $\triangle OAB$ is a right triangle at the origin.

 Label $D$ the point of tangency to the first, big circle, and label $E$ the point of tangency to the small circle.  $\angle D$ and $\angle E$ are both right angles as well.

 Label the center of the big circle $F$, and the small circle $G$.

 $\triangle DAF  \sim \triangle EAG \sim \triangle OAB$ because they each have one right angle, and also have a common angle $A$.

 The y-intercept is the length $OB$.

 We know that $DF = 2$ and $EG = 1$ because they are radii of circles.  From similarity, $\frac{DF}{EG} = \frac{FA}{GA}$.  Thus, $FA = 2GA$.

 Looking at line $FGA$, we know that $FG + GA = FA$.

 Thus, $FG + GA = 2GA$, meaning $FG = GA$.  Since $FG = 3$ because it's the distance of the centers of the circles, so $GA = 3$ as well.  This gives point $A$ as $(8,0)$.

 Since $\triangle EAG$ is a right triangle with $GA = 3$ and $GE = 1$, we know that $AE = \sqrt{GA^2 - GE^2} = \sqrt{8}$

 Thus, all triangles are $1:\sqrt{8}:3$ triangles.

 $\triangle OAC$ has as its middle side $OA$, which is $8$.  Thus, all sides are scaled up by a factor of $\sqrt{8}$, and $OB$, the shortest side, is $\sqrt{8}$.  This means the y-intercept is $\sqrt{8} = 2\sqrt{2}$, and the answer is $\boxed{D}$.

 



\textbf{Solution 2}

Because all the numbers are pretty small, and the correct answer choice is far enough apart from the others, just draw the diagram to scale with a ruler, see the y-intercept is a little less than 3.5, and get $\boxed{D}$.

 


	\item (\textit{AHSME 1996}) The midpoints of the sides of a regular hexagon $ABCDEF$ are joined to form a smaller hexagon. What fraction of the area of $ABCDEF$ is enclosed by the smaller hexagon? 


 
                    \begin{figure}[H]
                        \centering
                        \includegraphics[width=0.5\textwidth]{Problem/AHSME_I_1996_19_Problem_Figure_1.png}
                    \end{figure}
                    
$\text{(A)}\ \frac{1}{2}\qquad\text{(B)}\ \frac{\sqrt 3}{3}\qquad\text{(C)}\ \frac{2}{3}\qquad\text{(D)}\ \frac{3}{4}\qquad\text{(E)}\ \frac{\sqrt 3}{2}$

 \textbf{Answer: }D



\textbf{Solution}


                    \begin{figure}[H]
                        \centering
                        \includegraphics[width=0.5\textwidth]{Solution/AHSME_I_1996_19_Solution_Figure_1.png}
                    \end{figure}
                    
$\triangle OCD$ is copied six times to form the hexagon, so if we find the ratio of the area of the kite inside $\triangle OCD$ to the the area of $\triangle OCD$ itself, it will be the same ratio.

 Let $OC=CD=OD =2$ so that the area of the triangle is $\frac{\sqrt{3}s^2}{4} = \sqrt{3}$.

 Notice that $\triangle OCD$ is made up of a kite and two $30-60-90$ triangles.  The two hypotenuses of these two triangles form $CD$, so the hypotenuse of each triangle must be $\frac{2}{2} = 1$.  Thus, the legs of each triangle are $\frac{1}{2}$ and $\frac{\sqrt{3}}{2}$, and the area of two of these triangles is $\frac{1}{2}\cdot \frac{\sqrt{3}}{2} = \frac{\sqrt{3}}{4}$.

 Subtracting the area of the two triangles from the area of the equilateral triangle, we find that the area of the kite is $\sqrt{3} - \frac{\sqrt{3}}{4} = \frac{3\sqrt{3}}{4}$.

 Thus, the ratio of areas is $\frac{3}{4}$, which is option $\boxed{D}$.

 



\textbf{Solution 2}

We see six isosceles $120-30-30$ triangles at the vertices of the large hexagon.  If we let the side of the larger hexagon be $1$, we find that the two congruent sides of these triangles are both $\frac{1}{2}$.  By the \href{https://artofproblemsolving.com/wiki/index.php?title=Law\_of\_Sines}{Law of Sines}, we find that the third side $x$ is:

 $\frac{\frac{1}{2}}{\sin 30^\circ} = \frac{x}{\sin 120^\circ}$
$1 = \frac{x}{\frac{\sqrt{3}}{2}}$
$x = \frac{\sqrt{3}}{2}$

 This third side of the triangle is the length of the edge of the hexagon.  The ratio of the sides of the small hexagon to the large hexagon is $\frac{\sqrt{3}}{2}$, since the large hexagon has a unit side.  The ratio of the areas is the square of the ratio of the sides, which is $\frac{3}{4}$, or option $\boxed{D}$.

 


	\item (\textit{AHSME 1996}) In the xy-plane, what is the length of the shortest path from $(0,0)$ to $(12,16)$ that does not go inside the circle $(x-6)^{2}+(y-8)^{2}= 25$?


 $\text{(A)}\ 10\sqrt 3\qquad\text{(B)}\ 10\sqrt 5\qquad\text{(C)}\ 10\sqrt 3+\frac{ 5\pi}{3}\qquad\text{(D)}\ 40\frac{\sqrt{3}}{3}\qquad\text{(E)}\ 10+5\pi$

 \textbf{Answer: }C



\textbf{Solution}

The pathway from $A(0,0)$ to $D(12,16)$ will consist of three segments:

 1)  $\overline{AB}$, where $AB$ is tangent to the circle at point $B$.

 2)  $\overline{CD}$, where $CD$ is tangent to the circle at point $C$.

 3)  $\widehat {BC}$, where $BC$ is an arc around the circle.

 The actual path will go $A \rightarrow B \rightarrow C \rightarrow D$, so the actual segments will be in order $1, 3, 2$.

 Let $O$ be the center of the circle at $(6,8)$.

 $OA = 10$ and $OB = 5$ since $B$ is on the circle.  Since $\triangle OAB$ is a right triangle with right angle $B$, we find that $AB = \sqrt{10^2 - 5^2} = 5\sqrt{3}$.  This means that $\triangle OAB$ is a $30-60-90$ triangle with sides $5:5\sqrt{3}:10$.

 Thus, $\angle DOC = 60^\circ$, and $\angle AOB = 60^\circ$.  Since $\angle AOB$, $\angle DOC$ and $\angle BOC$ lie on a straight line, $\angle BOC$ must be $60^\circ$ as well.  Thus, the arc that we travel is a $60^\circ$ arc, and we travel $\frac{C}{6} = \frac{2\pi r}{6} = \frac{2\pi \cdot 5}{6} = \frac{5\pi}{3}$ around the circle.

 Thus, $AB = 5\sqrt{3}$, $\widehat {BC} = \frac{5\pi}{3}$, and ${CD} = 5\sqrt{3}$.  The total distance is $10\sqrt{3} + \frac{5\pi}{3}$, which is option $\boxed{C}$.

 



\textbf{Solution 2 (by elimination of options)}

The shortest path likely consists of two congruent line segments connected by an arc of the circle, so the total length can be expressed in the form
 \[a + b\pi,\]where \(a\) and \(b\) are positive real numbers.

 Using this observation, we can eliminate options A, B, and D, as they do not fit this form.

 Next, the total circumference of the circle is \(10\pi\). Since half the circumference is \(5\pi\), option E can be ruled out because it corresponds to exactly half the circle's circumference, which is unlikely to be the shortest path involving line segments.

 Therefore, by elimination of the other choices, the only remaining option is \(10\sqrt{3} + \frac{5\pi}{3}\), which matches option \(\boxed{C}\).

 


	\item (\textit{AHSME 1996}) Triangles $ABC$ and $ABD$ are isosceles with $AB=AC=BD$, and $BD$ intersects $AC$ at $E$. If $BD$ is perpendicular to $AC$, then $\angle C+\angle D$ is


 
                    \begin{figure}[H]
                        \centering
                        \includegraphics[width=0.5\textwidth]{Problem/AHSME_I_1996_21_Problem_Figure_1.png}
                    \end{figure}
                    
$\text{(A)}\ 115^\circ\qquad\text{(B)}\ 120^\circ\qquad\text{(C)}\ 130^\circ\qquad\text{(D)}\ 135^\circ\qquad\text{(E)}\ \text{not uniquely determined}$

 \textbf{Answer: }D



\textbf{Solution 1}

Redraw the figure as a concave pentagon $ADECB$:

 
                    \begin{figure}[H]
                        \centering
                        \includegraphics[width=0.5\textwidth]{Solution/AHSME_I_1996_21_Solution_Figure_1.png}
                    \end{figure}
                    

 The angles of the pentagon will still sum to $180^\circ \cdot 3 = 540^\circ$, regardless of whether the pentagon is concave or not.  As a quick proof, note that the nine angles of three original triangles $\triangle  AEB$, $\triangle CBE$, and $\triangle DEA$ all make up the angles of the pentagon without overlap.

 Since reflex $\angle E = 270^\circ$, we have:

 $\angle C + \angle B + \angle A + \angle D = 540^\circ - 270^\circ = 270^\circ$. 

 From isosceles $\triangle ABC$, we get $\angle B = \angle C$, so:

 $2\angle C + \angle D + \angle A = 270^\circ$

 From isosceles $\triangle ABD$, we get $\angle A = \angle D$, so:

 $2\angle C + 2\angle D = 270^\circ$
$\angle C + \angle D = 135^\circ$, which is answer $\boxed{D}$




\textbf{Solution 2}

Let $m\angle ABC = x$. By the isosceles triangle theorem, we have $m\angle ACB = x$ and $m\angle BAD = m\angle BDA$. Because the angles of a triangle sum to $\pi$, we have $m\angle BAC = \pi - 2x$, then $m\angle ABD=\frac{\pi}{2}-(\pi-2x)=2x-\frac{\pi}{2}$. Then we have $m\angle BAD + m\angle BDA = \pi - m\angle ABD$. Substituting, this becomes $2m\angle BDA = \pi-\left(\frac{\pi}{2}-(\pi-2x)\right) \to m \angle BDA = \frac{3 \pi}{4} - x$. Adding $x$, which is $m\angle ACB$, we have $m\angle BDA + m\angle ACB = \frac{3 \pi}{4} = 135^\circ \to \boxed{\textbf{D}}$



	\item (\textit{AHSME 1996}) Four distinct points, $A$, $B$, $C$, and $D$, are to be selected from $1996$ points 
evenly spaced around a circle. All quadruples are equally likely to be chosen. 
What is the probability that the chord $AB$ intersects the chord $CD$? 


 $\text{(A)}\ \frac{1}{4}\qquad\text{(B)}\ \frac{1}{3}\qquad\text{(C)}\ \frac{1}{2}\qquad\text{(D)}\ \frac{2}{3}\qquad\text{(E)}\ \frac{3}{4}$

 \textbf{Answer: }B



\textbf{Solution}

Let $WXYZ$ be a convex \href{https://artofproblemsolving.com/wiki/index.php?title=Cyclic\_quadrilateral}{cyclic quadrilateral} inscribed in a circle.  There are $\frac{\binom{4}{2}}{2} = 3$ ways to divide the points into two groups of two.

 If you pick $WX$ and $YZ$, you have two sides of the quadrilateral, which do not intersect.

 If you pick $XY$ and $ZW$, you have the other two sides of the quadrilateral, which do not intersect.

 If you pick $WY$ and $XZ$, you have the diagonals of the quadrilateral, which do intersect.

 Any four points on the original circle of $1996$ can be connected to form such a convex quadrilateral $WXYZ$, and only placing $A$ and $C$ as one of the diagonals of the figure will form intersecting chords.  Thus, the answer is $\frac{1}{3}$, which is option $\boxed{\text{B}}$.

 Notice that $1996$ is irrelevant to the solution of the problem; in fact, you may pick points from the entire \href{https://artofproblemsolving.com/wiki/index.php?title=Circumference}{circumference} of the circle.

 


	\item (\textit{AHSME 1996}) The sum of the lengths of the twelve edges of a rectangular box is $140$, and 
the distance from one corner of the box to the farthest corner is $21$. The total 
surface area of the box is 


 $\text{(A)}\ 776\qquad\text{(B)}\ 784\qquad\text{(C)}\ 798\qquad\text{(D)}\ 800\qquad\text{(E)}\ 812$

 \textbf{Answer: }B



\textbf{Solution}

Let $x, y$, and $z$ be the unique lengths of the edges of the box.  Each box has $4$ edges of each length, so:
 \[4x + 4y + 4z = 140 \ \Longrightarrow \ x + y + z = 35.\]The spacial diagonal (longest distance) is given by $\sqrt{x^2 + y^2 + z^2}$.  Thus, we have $\sqrt{x^2 + y^2 + z^2} = 21$, so $x^2 + y^2 + z^2 = 21^2$.

 Our target expression is the surface area of the box:

 
 \[S = 2xy + 2xz + 2yz.\]

 Since $S$ is a \href{https://artofproblemsolving.com/wiki/index.php?title=Elementary\_symmetric\_sum}{symmetric polynomial} of degree $2$, we try squaring the first equation to get:

 
 \[35^2 = (x + y + z)^2 = x^2 + y^2 + z^2 + 2xy +2yz + 2xz = 35^2.\]

 Substituting in our long diagonal and surface area expressions, we get: $21^2 + S = 35^2$, so $S = (35 + 21)(35 - 21) = 56\cdot 14 = 784$, which is option $\boxed{(\text{B})}$.

 


	\item (\textit{AHSME 1996}) The sequence $1, 2, 1, 2, 2, 1, 2, 2, 2, 1, 2, 2, 2, 2, 1, 2, 2, 2, 2, 2, 1, 2,\ldots$ consists of $1$’s separated by blocks of $2$’s with $n$ $2$’s in the $n^{th}$ block. The sum of the first $1234$ terms of this sequence is


 $\text{(A)}\ 1996\qquad\text{(B)}\ 2419\qquad\text{(C)}\ 2429\qquad\text{(D)}\ 2439\qquad\text{(E)}\ 2449$

 \textbf{Answer: }B



\textbf{Solution}

The sum of the first $1$ numbers is $1$

 The sum of the next $2$ numbers is $2 + 1$

 The sum of the next $3$ numbers is $2 + 2 + 1$

 In general, we can write "the sum of the next $n$ numbers is $1 + 2(n-1)$", where the word "next" follows the pattern established above.

 Thus, we first want to find what triangular numbers $1234$ is between.  By plugging in various values of $n$ into $f(n) = \frac{n(n+1)}{2}$, we find:

 $f(50) = 1275$
$f(49) = 1225$

 Thus, we want to add up all those sums from "next $1$ number" to the "next $49$ numbers", which will give us all the numbers up to and including the $1225^{th}$ number.  Then, we can manually tack on the remaining $2$s to hit $1234$.

 We want to find:

 $\sum_{n=1}^{49} 1 + 2(n-1)$
$\sum_{n=1}^{49} 2n - 1$
$\sum_{n=1}^{49} 2n - \sum_{n=1}^{49} 1$
$2 \sum_{n=1}^{49} n - 49$
$2\cdot \frac{49\cdot 50}{2} - 49$
$49^2$
$2401$

 Thus, the sum of the first $1225$ terms is $2401$.  We have to add $9$ more $2$s to get to the $1234^{th}$ term, which gives us $2419$, or option $\boxed{B}$.

 Note:  If you notice that the above sums form $1 + 3 + 5 + 7... + (2n-1) = n^2$, the fact that $49^2$ appears at the end should come as no surprise.

 



\textbf{Solution 2 (Alcumus)}

The $k$th appearance of 1 is at position $1 + 2 + \dots + k = \frac{k(k + 1)}{2}$. Then there are $k$ 1's and $\frac{k(k + 1)}{2} - k = \frac{k(k - 1)}{2}$ 2's among the first $\frac{k(k + 1)}{2}$ numbers, so the sum of these $\frac{k(k + 1)}{2}$ terms is $k + k(k - 1) = k^2$.

 When $k = 49$, $\frac{k(k + 1)}{2} = 1225$, and when $k = 50$, $\frac{k(k + 1)}{2} = 1275$.

 The sum of the first 1225 terms is $49^2 = 2401$. The numbers in positions 1226 through 1234 are all 2's, so their sum is $(1234 - 1226 + 1) \cdot 2 = 18$. Therefore, the sum of the first 1234 terms is $2401 + 18 = \boxed{2419}$.

 


	\item (\textit{AHSME 1996}) Given that $x^2 + y^2 = 14x + 6y + 6$, what is the largest possible value that $3x + 4y$ can have? 


 $\text{(A)}\ 72\qquad\text{(B)}\ 73\qquad\text{(C)}\ 74\qquad\text{(D)}\ 75\qquad\text{(E)}\ 76$

 \textbf{Answer: }B



\textbf{Solution 1}

Complete the square to get 
 \[(x-7)^2 + (y-3)^2 = 64.\]Applying Cauchy-Schwarz directly,
 \[64\cdot25=(3^2+4^2)((x-7)^2 + (y-3)^2) \ge (3(x-7)+4(y-3))^2.\]
 \[40 \ge 3x+4y-33\]
 \[3x+4y \le 73.\]Thus our answer is $\boxed{(B)}$.

 



\textbf{Solution 2 (Geometric)}

The first equation is a \href{https://artofproblemsolving.com/wiki/index.php?title=Circle}{circle}, so we find its center and \href{https://artofproblemsolving.com/wiki/index.php?title=Radius}{radius} by \href{https://artofproblemsolving.com/wiki/index.php?title=Completing\_the\_square}{completing the square}: $x^2 - 14x + y^2 - 6y = 6$, so 
 \[(x-7)^2 + (y-3)^2 = (x^2- 14x + 49) + (y^2 - 6y + 9) = 6 + 49 + 9 = 64.\]

 So we have a circle centered at $(7,3)$ with radius $8$, and we want to find the max of $3x + 4y$.

 The set of lines $3x + 4y = A$ are all \href{https://artofproblemsolving.com/wiki/index.php?title=Parallel}{parallel}, with slope 

 


	\item (\textit{AHSME 1996}) An urn contains marbles of four colors: red, white, blue, and green. When four marbles are drawn without replacement, the following events are equally likely: 


 (a) the selection of four red marbles; 


 (b) the selection of one white and three red marbles; 


 (c) the selection of one white, one blue, and two red marbles; and 


 (d) the selection of one marble of each color. 


 What is the smallest number of marbles satisfying the given condition? 


 $\text{(A)}\ 19\qquad\text{(B)}\ 21\qquad\text{(C)}\ 46\qquad\text{(D)}\ 69\qquad\text{(E)}\ \text{more than 69}$

 \textbf{Answer: }B



\textbf{Solution}

Let the bag contain $n$ marbles total, with $r, w, b, g$ representing the number of red, white, blue, and green marbles, respectively.  Note that $r + w + b + g = n$.

 The number of ways to select four red marbles out of the set of marbles without replacement is:

 
 \[\binom{r}{4} = \frac{r!}{24\cdot (r -4)!}\]

 The number of ways to select one white and three red marbles is:

 
 \[\binom{w}{1}\binom{r}{3} = \frac{w\cdot r!}{6\cdot (r - 3)!}\] 

 The number of ways to select one white, one blue, and two red marbles is:

 
 \[\binom{w}{1}\binom{b}{1} \binom{r}{2} = \frac{wb\cdot r!}{2(r-2)!}\]

 The number of ways to select one marble of each colors is:

 
 \[\binom{w}{1}\binom{b}{1} \binom{g}{1}\binom{r}{1} = wbgr\]Setting the first and second statements equal, we find:

 
 \[\frac{r!}{24\cdot (r -4)!}  = \frac{w\cdot r!}{6\cdot (r - 3)!}\]

 \[r - 3  = 4w\]

 Setting the first and third statements equal, we find: 

 
 \[\frac{r!}{24\cdot (r -4)!} = \frac{wb\cdot r!}{2(r-2)!}\]

 \[(r-3)(r-2) = 12wb\]

 Setting the last two statements equal, we find:

 
 \[\frac{wb\cdot r!}{2(r-2)!} = wbgr\]

 \[r - 1 = 2g\]

 These are all the "linking equations" that are needed; the transitive property of equality makes the other three equalities unnecessary.

 From the first equation, we know that $r$ must be $3$ more than a multiple of $4$, or that $r \equiv 3 \mod 4$

 Putting the first equation into the second equation, we find $r-2 = 3b$.  Therefore, $r \equiv 2 \mod 3$.  Using the \href{https://artofproblemsolving.com/wiki/index.php?title=Chinese\_Remainder\_Theorem}{Chinese Remainder Theorem}, we find that $r \equiv 11 \mod 12$.

 The third equation gives no new restrictions on $r$; it is already odd by the first equation.  

 Thus, the minimal positive value of $r$ is $11$.  This requires $g=\frac{r - 1}{2} = 5$ by the third equation, and $w = \frac{r-3}{4} = 2$ by the first equation.  Finally, the second equation gives $b = \frac{(r-3)(r-2)}{12w} = 3$.

 The minimal total number of marbles is $11 + 5 + 2 + 3 = \boxed{21}$, which is option $\boxed{B}$.

 


	\item (\textit{AHSME 1996}) Consider two solid spherical balls, one centered at $(0, 0,\frac{21}{2})$ with radius $6$, and the other centered at $(0, 0, 1)$ with radius $\frac{9}{2}$. How many points with only integer coordinates (lattice points) are there in the intersection of the balls? 


 $\text{(A)}\ 7\qquad\text{(B)}\ 9\qquad\text{(C)}\ 11\qquad\text{(D)}\ 13\qquad\text{(E)}\ 15$

 \textbf{Answer: }D



\textbf{Solution 1}

The two equations of the balls are

 
 \[x^2 + y^2 + \left(z - \frac{21}{2}\right)^2 \le 36\]

 \[x^2 + y^2 + (z - 1)^2 \le \frac{81}{4}\]

 Note that along the $z$ axis, the first ball goes from $10.5 \pm 6$, and the second ball goes from $1 \pm 4.5$.  The only integer value that $z$ can be is $z=5$.

 Plugging that in to both equations, we get:

 
 \[x^2 + y^2 \le \frac{23}{4}\]

 \[x^2 + y^2 \le \frac{17}{4}\]

 The second inequality implies the first inequality, so the only condition that matters is the second inequality.

 From here, we do casework, noting that $|x|, |y| \le 3$:

 For $x=\pm 2$, we must have $y=0$.  This gives $2$ points.

 For $x = \pm 1$, we can have $y\in \{-1, 0, 1\}$.  This gives $2\cdot 3 = 6$ points.

 For $x = 0$, we can have $y \in \{-2, -1, 0, 1, 2\}$.  This gives $5$ points.

 Thus, there are $\boxed{13}$ possible points, giving answer $\boxed{D}$.

 



\textbf{Solution 2 (minimal 3D geometry)}

Because both spheres have their centers on the x-axis, we can simplify the graph a bit by looking at a 2-dimensional plane (the previous z-axis is the new x-axis while the y-axis remains the same).

 The spheres now become circles with centers at $(1,0)$ and $(\frac{21}{2},0)$.  They have radii $\frac{9}{2}$ and $6$, respectively.Let circle $A$ be the circle centered on $(1,0)$ and circle $B$ be the one centered on $(\frac{21}{2},0)$.

 The point on circle $A$ closest to the center of circle $B$ is $(\frac{11}{2},0)$.  The point on circle B closest to the center of circle $A$ is $(\frac{11}{2},0)$.

 Taking a look back at the 3-dimensional coordinate grid with the spheres, we can see that their intersection appears to be a circle with congruent "dome" shapes on either end.  Because the tops of the "domes" are at $(0,0,\frac{9}{2})$ and $(0,0,\frac{11}{2})$, respectively, the lattice points inside the area of intersection must have z-value $5$ (because $5$ is the only integer between $\frac{9}{2}$ and $\frac{11}{2}$).  Thus, the lattice points in the area of intersection must all be on the 2-dimensional circle.  The radius of the circle will be the distance from the z-axis.

 Now, looking at the 2-dimensional coordinate plane, we see that the radius of the circle (now the distance from the x-axis, because there is no more z-axis) is the altitude of a triangle with two points on centers of circles $A$ and $B$ and third point at the first quadrant intersection of the circles.  Let's call that altitude $h$.

 We know all three side lengths of this triangle:  $\frac{9}{2}$ (the radius of circle $A$), $6$ (the radius of circle $B$), and $\frac{19}{2}$ (the distance between the centers of circles $A$ and $B$).  We can now find the area of the triangle using Heron's formula:

 
 \[s=\frac{\frac{9}{2}+6+\frac{19}{2}}{2}=20\]

 \[A=\sqrt{(20)(20-\frac{9}{2})(20-6)(20-\frac{19}{2})}=\sqrt{110}\]

 Using the area of the triangle, we can find that altitude $h$ from the x-axis:

 
 \[\sqrt{110}=\frac{h*\frac{19}{2}}{2}\]

 \[h=\frac{4*\sqrt{110}}{19}\]

 Remember, the altitude $h$ is also the radius of the circle containing all the solutions to the problem.

 Going back to the 3-dimensional grid and looking at the circle, we can again make the figure 2-dimensional.

 The radius $h$ of the circle in a 2-dimensional plane is $\frac{4\sqrt{110}}{19}$, a little greater than $2$.  We know that $h>2$ because $4\sqrt{110}>4*10$, and $\frac{40}{19}>2$.

 Finally, looking at a circle with radius slightly larger than $2$, we see that there are $\boxed{13}$ lattice points within it $\boxed{D}$.

 



\textbf{Solution 3}

Note that the spheres are on the same $x$ and $y$ axis. Therefore, we can draw the spheres so that only the $x$  and $z$ axis are featured. 

 
                    \begin{figure}[H]
                        \centering
                        \includegraphics[width=0.5\textwidth]{Solution/AHSME_I_1996_27_Solution_Figure_1.jpg}
                    \end{figure}
                    $A = (1, 0, 5)$. $\sqrt{1^2+4^2} = \sqrt{17}<\frac92$, $A$ is inside the smaller sphere. $\sqrt{1^2+(\frac{11}{2})^2} = \frac{\sqrt{126}}{2}<6$, $A$ is inside the larger sphere. Point $A$ is inside the intersection.

 $B = (2, 0, 5)$, $\sqrt{2^2+4^2} = \sqrt{20}<\frac92$, $B$ is inside the smaller sphere. $\sqrt{2^2+(\frac{11}{2})^2} = \frac{\sqrt{137}}{2}<6$, $B$ is inside the larger sphere. Point $B$ is inside the intersection.

 $C = (3, 0, 5)$, $\sqrt{3^2+4^2} = \sqrt{25}>\frac92$, $C$ is outside the smaller sphere. $\sqrt{3^2+(\frac{11}{2})^2} = \frac{\sqrt{157}}{2}>6$, $C$ is outside the larger sphere. Point $C$ is outside the intersection.

 The intersection of $2$ spheres is a circle, by drawing the circle flat we can see that there are $4$ more points within the intersection. 

 We can see that points with the same $x$-coordinates with $A$ are still inside the circle. $B$ is very close to the circle, the points with the same $x$-coordinates with $B$ are not necessarily inside the circle.  

 For example, $D = (2,1,5)$, $\sqrt{2^2+1^2+4^2}=\sqrt{21}>\frac92$, $D$ is outside the smaller sphere. $\sqrt{2^2 + 1^2 + (\frac{11}{2})^2} = \frac{\sqrt{141}}{2}<6$, $D$ is inside the larger sphere. Point $D$ is outside the intersection.

 
                    \begin{figure}[H]
                        \centering
                        \includegraphics[width=0.5\textwidth]{Solution/AHSME_I_1996_27_Solution_Figure_2.jpg}
                    \end{figure}
                    Therefore, the answer is $9+4 = \boxed{\textbf{(D) } 13}$.

 


	\item (\textit{AHSME 1996}) On a $4\times 4\times 3$ rectangular parallelepiped, vertices $A$, $B$, and $C$ are adjacent to vertex $D$. The perpendicular distance from $D$ to the plane containing
$A$, $B$, and $C$ is closest to 


 
                    \begin{figure}[H]
                        \centering
                        \includegraphics[width=0.5\textwidth]{Problem/AHSME_I_1996_28_Problem_Figure_1.png}
                    \end{figure}
                    
$\text{(A)}\ 1.6\qquad\text{(B)}\ 1.9\qquad\text{(C)}\ 2.1\qquad\text{(D)}\ 2.7\qquad\text{(E)}\ 2.9$

 \textbf{Answer: }C



\textbf{Solution 1}

By placing the cube in a coordinate system such that $D$ is at the origin, $A(0,0,3)$, $B(4,0,0)$, and $C(0,4,0)$, we find that the equation of plane $ABC$ is:

 
 \[\frac{x}{4} + \frac{y}{4} + \frac{z}{3} = 1,\] so $3x + 3y + 4z - 12 = 0.$ The equation for the distance of a point $(a,b,c)$ to a plane $Ax + By + Cz + D = 0$ is given by:

 
 \[\frac{|Aa + Bb + Cc + D|}{\sqrt{A^2 + B^2 + C^2}}.\]

 Note that the capital letters are coefficients, while the lower case is the point itself. Thus, the distance from the origin (where $a=b=c=0$) to the plane is given by:

 
 \[\frac{D}{\sqrt{A^2 + B^2 + C^2}} = \frac{12}{\sqrt{9 + 9 + 16}} = \frac{12}{\sqrt{34}}.\]

 Since $\sqrt{34} < 6$, this number should be just a little over $2$, and the correct answer is $\boxed{\text{(C)}}$.

 Note that the equation above for the distance from a point to a plane is a 3D analogue of the 2D case of the \href{https://artofproblemsolving.com/wiki/index.php?title=Distance\_formula}{distance formula}, where you take the distance from a point to a line.  In the 2D case, both $c$ and $C$ are set equal to $0$.

 



\textbf{Solution 2}

Let $x$ be the desired distance. Recall that the volume of a pyramid is given by $\frac{1}{3}\cdot h \cdot B$, where $B$ is the area of the base and $h$ is the height. Consider pyramid $ABCD$. Letting $ABC$ be the base, the volume of $ABCD$ is given by $\frac{1}{3} \cdot x \cdot [ABC]$, but if we let $BCD$ be the base, the volume is given by $\frac{1}{3} \cdot [BCD]\cdot [AD] = \frac{1}{3} \cdot [\frac{1}{2} \cdot 4 \cdot 4] \cdot 3 = 8$. Clearly, these two volumes must be equal, so we get the equation $\frac{1}{3}\cdot x \cdot[ABC]=8$. Thus, to find $x$, we just need to find $[ABC]$.

 By the Pythagorean Theorem, $AB=\sqrt{AD^2+DB^2}=5$, $AC=\sqrt{AD^2+DC^2}=5$, $BC=\sqrt{BD^2+DC^2}=4\sqrt{2}$. 

 The altitude to $BC$ in triangle $ABC$ has length $\sqrt{AC^2-\frac{BC}{2}^2}=\sqrt{17}$, so $[ABC]=\frac{1}{2}\cdot 4\sqrt{2} \cdot \sqrt{17} = 2\sqrt{34}$. Then $x=\frac{24}{[ABC]}=\frac{24}{2\sqrt{34}}=\frac{6\sqrt{34}}{17}$ or about $2.1$. The answer is $\boxed{C}$.

 


	\item (\textit{AHSME 1996}) If $n$ is a positive integer such that $2n$ has $28$ positive divisors and $3n$ has $30$ positive divisors, then how many positive divisors does $6n$ have?


 $\text{(A)}\ 32\qquad\text{(B)}\ 34\qquad\text{(C)}\ 35\qquad\text{(D)}\ 36\qquad\text{(E)}\ 38$

 \textbf{Answer: }C



\textbf{Solution 1}

Working with the second part of the problem first, we know that $3n$ has $30$ divisors.  We try to find the various possible prime factorizations of $3n$ by splitting $30$ into various products of $1, 2$ or $3$ integers.

 $30 \rightarrow p^{29}$
$2 \cdot 15 \rightarrow pq^{14}$
$3\cdot 10 \rightarrow p^2q^9$
$5\cdot 6 \rightarrow p^4q^5$
$2\cdot 3\cdot 5 \rightarrow pq^2r^4$

 The variables $p, q, r$ are different prime factors, and one of them must be $3$.  We now try to count the factors of $2n$, to see which prime factorization is correct and has $28$ factors.

 In the first case, $p=3$ is the only possibility.  This gives $2n = 2\cdot p^{28}$, which has $2\cdot {29}$ factors, which is way too many.

 In the second case, $p=3$ gives $2n = 2q^{14}$.  If 

 


	\item (\textit{AHSME 1996}) A hexagon inscribed in a circle has three consecutive sides each of length $3$ and three consecutive sides each of length $5$. The chord of the circle that divides the hexagon into two trapezoids, one with three sides each of length $3$ and the other with three sides each of length $5$, has length equal to $\frac{m}{n}$, where $m$ and $n$ are relatively prime positive integers. Find $m+n$.


 $\text{(A)}\ 309\qquad\text{(B)}\ 349\qquad\text{(C)}\ 369\qquad\text{(D)}\ 389\qquad\text{(E)}\ 409$

 \textbf{Answer: }E



\textbf{Solution 1 (Alcumus)}

In hexagon $ABCDEF$, let $AB=BC=CD=3$ and let $DE=EF=FA=5$. Since arc $BAF$ is one third of the circumference of the circle, it follows that $\angle BCF = \angle BEF=60^{\circ}$. Similarly, $\angle CBE =\angle CFE=60^{\circ}$. Let $P$ be the intersection of $\overline{BE}$ and $\overline{CF}$, $Q$ that of $\overline{BE}$ and $\overline{AD}$, and $R$ that of $\overline{CF}$ and $\overline{AD}$. Triangles $EFP$ and  $BCP$ are equilateral, and by symmetry, triangle $PQR$ is isosceles and thus also equilateral. 
                    \begin{figure}[H]
                        \centering
                        \includegraphics[width=0.5\textwidth]{Solution/AHSME_I_1996_30_Solution_Figure_1.png}
                    \end{figure}
                    

 Furthermore, $\angle BAD$ and $\angle BED$ subtend the same arc, as do $\angle ABE$ and $\angle ADE$. Hence triangles $ABQ$ and $EDQ$ are similar. Therefore, 
 \[\frac{AQ}{EQ}=\frac{BQ}{DQ}=\frac{AB}{ED}=\frac{3}{5}.\] It follows that 
 \[\frac{\frac{AD-PQ}{2}}{PQ+5} =\frac{3}{5}\quad \mbox {and}\quad \frac{3-PQ}{\frac{AD+PQ}{2}}=\frac{3}{5}.\] Solving the two equations simultaneously yields $AD=360/49,$ so $m+n=\boxed{409}$.

 



\textbf{Solution 2}

All angle measures are in degrees.Let the first trapezoid be $ABCD$, where $AB=BC=CD=3$.  Then the second trapezoid is $AFED$, where $AF=FE=ED=5$.  We look for $AD$.

 Since $ABCD$ is an isosceles trapezoid, we know that $\angle BAD=\angle CDA$ and, since $AB=BC$, if we drew $AC$, we would see $\angle BCA=\angle BAC$.  Anyway, $\widehat{AB}=\widehat{BC}=\widehat{CD}$ ($\widehat{AB}$ means arc AB).  Using similar reasoning, $\widehat{AF}=\widehat{FE}=\widehat{ED}$.

 Let $\widehat{AB}=2\phi$ and $\widehat{AF}=2\theta$.  Since $6\theta+6\phi=360$ (add up the angles), $2\theta+2\phi=120$ and thus $\widehat{AB}+\widehat{AF}=\widehat{BF}=120$.  Therefore, $\angle FAB=\frac{1}{2}\widehat{BDF}=\frac{1}{2}(240)=120$.  $\angle CDE=120$ as well.

 Now I focus on triangle $FAB$.  By the Law of Cosines, $BF^2=3^2+5^2-30\cos{120}=9+25+15=49$, so $BF=7$.  Seeing $\angle ABF=\theta$ and $\angle AFB=\phi$, we can now use the Law of Sines to get:
 \[\sin{\phi}=\frac{3\sqrt{3}}{14}\;\text{and}\;\sin{\theta}=\frac{5\sqrt{3}}{14}.\]

 Now I focus on triangle $AFD$.  $\angle AFD=3\phi$ and $\angle ADF=\theta$, and we are given that $AF=5$, so
 \[\frac{\sin{\theta}}{5}=\frac{\sin{3\phi}}{AD}.\]We know $\sin{\theta}=\frac{5\sqrt{3}}{14}$, but we need to find $\sin{3\phi}$.  Using various identities, we see
 \begin{align*}\sin{3\phi}&=\sin{(\phi+2\phi)}=\sin{\phi}\cos{2\phi}+\cos{\phi}\sin{2\phi}\\ &=\sin{\phi}(1-2\sin^2{\phi})+2\sin{\phi}\cos^2{\phi}\\ &=\sin{\phi}\left(1-2\sin^2{\phi}+2(1-\sin^2{\phi})\right)\\ &=\sin{\phi}(3-4\sin^2{\phi})\\ &=\frac{3\sqrt{3}}{14}\left(3-\frac{27}{49}\right)=\frac{3\sqrt{3}}{14}\left(\frac{120}{49}\right)=\frac{180\sqrt{3}}{343} \end{align*}
Returning to finding $AD$, we remember 
 \[\frac{\sin{\theta}}{5}=\frac{\sin{3\phi}}{AD}\;\text{so}\;AD=\frac{5\sin{3\phi}}{\sin{\theta}}.\]Plugging in and solving, we see $AD=\frac{360}{49}$. Thus, the answer is $360 + 49 = 409$, which is answer choice $\boxed{\textbf{(E)}}$.

 



\textbf{Solution 3}

Let $x$ be the desired length. One can use Parameshvara's circumradius formula, which states that for a cyclic quadrilateral with sides $a, b, c, d$ the circumradius $R$ satisfies 
 \[R^2=\frac{1}{16}\cdot\frac{(ab+cd)(ac+bd)(ad+bc)}{(s-a)(s-b)(s-c)(s-d)},\] where $s$ is the semiperimeter. Applying this to the trapezoid with sides $3, 3, 3, x$, we see that many terms cancel and we are left with 
 \[R^2=\frac{27}{9-x}\] Similar canceling occurs for the trapezoid with sides $5, 5, 5, x$, and since the two quadrilaterals share the same circumradius, we can equate: 
 \[\frac{27}{9-x}=\frac{125}{15-x}\] Solving for $x$ gives $x=\frac{360}{49}$, so the answer is $\fbox{(E) 409}$.

 



\textbf{Solution 4}


                    \begin{figure}[H]
                        \centering
                        \includegraphics[width=0.5\textwidth]{Solution/AHSME_I_1996_30_Solution_Figure_2.png}
                    \end{figure}
                    Note that minor arc $\overarc{AB}$ is a third of the circumference, therefore, $\angle AOB = 120^{\circ}$. Major arc $\overarc{AB}$$=240^{\circ}$, $\angle ACB = 120^{\circ}$

 By the Law of Cosine, $AB = \sqrt{ 3^2 + 5^2 - 2 \cdot 3 \cdot 5 \cdot \cos 120^{\circ}} = \sqrt{ 9 + 25 + 15 } = 7$
$\frac{\angle AOB}{2} = 60^{\circ}$, therefore, $r = \frac{\frac{AB}{2}}{\sin 60^{\circ}} = \frac{\frac{7}{2}}{ \frac{\sqrt{3}}{2} } = \frac{7\sqrt{3}}{3}$
$\sin \frac{\theta}{2} = \frac{\frac32}{r} = \frac{3}{2r} = \frac{3}{2 \cdot \frac{7\sqrt{3}}{3}} = \frac{3 \sqrt{3}}{14}$

 Let $x$ be the length of the chord, $\sin \frac{3 \theta}{2} = \frac{\frac{x}{2}}{r}$

 By the triple angle formula, $\sin \frac{3 \theta}{2} = 3 \cdot \sin \frac{\theta}{2} - 4 \cdot \sin(\frac{ \theta}{2})^3 = 3 \cdot \frac{3 \sqrt{3}}{14} - 4 \cdot (\frac{3 \sqrt{3}}{14})^3$
$x = 2 \cdot \frac{7\sqrt{3}}{3} \cdot [3 \cdot \frac{3 \sqrt{3}}{14} - 4 \cdot (\frac{3 \sqrt{3}}{14})^3] = 2 \cdot \frac{7\sqrt{3}}{3} \cdot (\frac{9\sqrt{3}}{14} - \frac{82\sqrt{3}}{2 \cdot 7^3}) = 9 - \frac{81}{49} = \frac{360}{49}$

 Therefore, the answer is $\boxed{\textbf{(E) } 409}$.

 



\textbf{Solution 5}


                    \begin{figure}[H]
                        \centering
                        \includegraphics[width=0.5\textwidth]{Solution/AHSME_I_1996_30_Solution_Figure_3.png}
                    \end{figure}
                    Note that minor arc $\overarc{AB}$ is a third of the circumference, therefore, $\angle AOB = 120^{\circ}$.

 $\sin \frac{\alpha}{2} = \frac{\frac32}{r}$, $\sin \frac{\alpha}{2} = \frac{3}{2r}$
$\sin \frac{120^{\circ}-\alpha}{2} = \frac{\frac52}{r}$, $\sin (60^{\circ} - \frac{\alpha}{2}) = \frac{5}{2r}$
$\frac{\sin \frac{\alpha}{2}}{\sin (60^{\circ} - \frac{\alpha}{2}) } = \frac{\frac{3}{2r}}{\frac{5}{2r}} = \frac35$, $5 \cdot \sin \frac{\alpha}{2} = 3 \cdot \sin (60^{\circ} - \frac{\alpha}{2})$
$5 \cdot \sin \frac{\alpha}{2} = 3 ( \sin 60^{\circ} \cos \frac{\alpha}{2} - \sin \frac{\alpha}{2} \cos 60^{\circ}) = 3 ( \frac{\sqrt{3}}{2} \cdot \cos \frac{\alpha}{2} - \frac12 \cdot \sin \frac{\alpha}{2})$
$13 \cdot \sin \frac{\alpha}{2} = \frac{3\sqrt{3}}{2} \cdot \cos \frac{\alpha}{2}$

 Let $\sin \frac{\alpha}{2} = a$, $\cos \frac{\alpha}{2} = \sqrt{1-a^2}$, $13a = \frac{3\sqrt{3}}{2} \cdot  \sqrt{1-a^2}$
$169a^2 = 27-27a^2$, $196a^2=27$, $\sin \frac{\alpha}{2} = a = \sqrt{\frac{27}{196}} = \frac{3 \sqrt{3}}{14}$

 Let $x$ be the length of the chord, $\sin \frac{3 \alpha}{2} = \frac{\frac{x}{2}}{r}$

 By the triple angle formula, $\sin \frac{3 \alpha}{2} = 3 \cdot \sin \frac{\alpha}{2} - 4 \cdot \sin(\frac{ \alpha}{2})^3 = 3 \cdot \frac{3 \sqrt{3}}{14} - 4 \cdot (\frac{3 \sqrt{3}}{14})^3$
$x = 2 \cdot \frac{7\sqrt{3}}{3} \cdot [3 \cdot \frac{3 \sqrt{3}}{14} - 4 \cdot (\frac{3 \sqrt{3}}{14})^3] = 2 \cdot \frac{7\sqrt{3}}{3} \cdot (\frac{9\sqrt{3}}{14} - \frac{82\sqrt{3}}{2 \cdot 7^3}) = 9 - \frac{81}{49} = \frac{360}{49}$

 Therefore, the answer is $\boxed{\textbf{(E) } 409}$.

 



\textbf{Solution 6 (Ptolemy's theorem)}


                    \begin{figure}[H]
                        \centering
                        \includegraphics[width=0.5\textwidth]{Solution/AHSME_I_1996_30_Solution_Figure_4.png}
                    \end{figure}
                    Note that major arc $\overarc{AE}$ is two thirds of the circumference, therefore, $\angle AFE = 120^{\circ}$.

 By the Law of Cosine, $AE= \sqrt{ 3^2 + 5^2 - 2 \cdot 3 \cdot 5 \cdot \cos 120^{\circ}} = \sqrt{ 9 + 25 + 15 } = 7$

 By the Ptolemy's theorem of quadrilateral $ABDE$, $AD \cdot BE = AB \cdot DE + BD \cdot AE$, $AD = BE$, $AD^2= 3 \cdot 5 + 7^2 = 64$, $AD = 8$

 By the Ptolemy's theorem of quadrilateral $ABCD$, $AC \cdot BD = BC \cdot AD + AB \cdot CD$, $7AC = 3 \cdot 8 + 3 \cdot 5 = 39$, $AC = \frac{39}{7}$

 By the Ptolemy's theorem of quadrilateral $ABCF$, $AC \cdot BF = AB \cdot CF + BC \cdot AF$, $AC = BF$, $(\frac{39}{7})^2 = 3 \cdot CF + 3 \cdot 3$, $CF = \frac{360}{49}$

 Therefore, the answer is $\boxed{\textbf{(E) } 409}$.

 


\end{enumerate}